\section{The global isometry and the measured boundary}
\label{global:section}

We now apply Theorem~\ref{local:extension} to the data constructed in
Section~\ref{data:section}. The source evaluations identify which
points should correspond. Their separation makes this correspondence
single-valued, and their behavior at the boundary prevents its escape
from the interior. The local theorem then removes every interior
frontier of the isometric correspondence.
The use of a maximal matching set and limits of source evaluations
follows the source-embedding architecture in
\citep[Section~6.1]{LLS20}; Theorem~\ref{local:extension} supplies the
smooth continuation step needed here.

Retain the extended manifolds $M_j^e$, their common cap $E$, the open
source set $U_0\Subset E$, and the functionals $I_j$ of
Equation~\eqref{data:evaluations}. For this section an \emph{admissible
match} is a smooth local isometry
\[
 \Phi:W\longrightarrow (M_1^e)^\circ,
 \qquad W\subset(M_2^e)^\circ,
\]
where $W$ is connected, open, and contains $E$, such that
\begin{equation}\label{global:matching}
 \Phi|_E=\Id,\qquad I_2(p)=I_1(\Phi(p))\quad(p\in W).
\end{equation}
Thus $\Phi^*g_1=g_2$ on $W$. The identity on $E$ is an admissible
match by Proposition~\ref{data:cap}.

\begin{proposition}[No interior frontier]\label{global:frontier}
There is an isometry
\[
 \Phi:(M_2^e)^\circ\longrightarrow(M_1^e)^\circ
\]
onto the first interior satisfying Equation~\eqref{global:matching}.
\end{proposition}

\begin{proof}
\emph{The maximal match and its Green kernel.}\quad
Two admissible matches agree wherever their domains overlap, since
$I_1$ separates interior points by Lemma~\ref{data:jet-family}.
Their union therefore defines a single smooth local isometry on
the union $W$ of all admissible domains. This union is connected
because every domain contains the connected set $E$. The union
map is injective: equal images give equal $I_2$ evaluations, hence
equal original points. It is itself an admissible match and is
maximal by construction.

We first record the Green equality on this larger region:
\begin{equation}\label{global:green-equality}
 G_2(p,q)=G_1(\Phi(p),\Phi(q)),
             \qquad p,q\in W,\quad p\ne q.
\end{equation}
For fixed $p\in W$, equality of its evaluations and the common
source density give this equality for the other variable in $U_0$,
initially in the distributional sense. Both Green functions are
smooth away from their pole. After pullback by the local isometry,
their difference is harmonic on $W\setminus\{p\}$. This set is
connected, and equality holds on a nonempty open part of it.
Unique continuation proves Equation~\eqref{global:green-equality}.

\emph{A frontier partner and common coordinates.}\quad
Suppose $W$ is not the entire second interior. Connectedness gives
an interior frontier point $p$. If $p_k\in W$ tends to $p$,
compactness of $M_1^e$ gives a subsequential limit $q$ of
$\Phi(p_k)$. For every fixed source $f$, continuity up to the
boundary gives
\[
 I_2(p)(f)=I_1(q)(f).
\]
The point $q$ is interior: all boundary evaluations are zero,
whereas $I_2(p)$ is nonzero. Any two such limits have the same
evaluations and so coincide. It follows that every matched sequence
tending to $p$ has partners tending to this unique $q$.
In particular, nearby matched points have partners in any prescribed
neighborhood of $q$.

Both $p$ and $q$ are outside $\overline{U_0}$. The first statement
holds because $\overline{U_0}\subset E\subset W$. If $q$ belonged
to $\overline{U_0}$, the common-cap identification would give
$I_1(q)=I_2(q)$; separation on the second manifold would force
$p=q\in E$, a contradiction.

Choose sources $f_1,\ldots,f_n$ so that
\[
 y=(u_{f_1}^2,\ldots,u_{f_n}^2)
\]
has independent differentials at $p$, and put
$x=(u_{f_1}^1,\ldots,u_{f_n}^1)$ near $q$.
The second tuple is also a coordinate system. Indeed the isometries
at the points $p_k$ give equality of the gradient Gram matrices
\[
 \bigl(\langle du_{f_a}^2,du_{f_b}^2\rangle_{g_2}(p_k)\bigr)_{a,b}
 =
 \bigl(\langle du_{f_a}^1,du_{f_b}^1\rangle_{g_1}(\Phi(p_k))\bigr)_{a,b}.
\]
The first limit is positive definite, and so is the second.
The coordinate values at $p,q$ coincide. Restricting the charts
gives a common coordinate ball $V$ about this value, with closures
away from $\overline{U_0}$. The preceding convergence observation
ensures that every already matched point sufficiently near $p$
has its partner in the first chart. In these coordinates its
identification is $x=y$.

Let $A\subset V$ be the coordinate image of the already matched
second-side points. It is open and has the center value in its
relative frontier. The two coordinate metrics and every paired
source function agree on $A$. By Lemma~\ref{data:jet-family},
choose finitely many first-side source functions whose ordinary
Hessians span the hyperplane annihilated by $g_1^{-1}$ at the
center. After shrinking $V$, their spanning rank has a positive
lower bound throughout $V$. Their second-side partners are
harmonic there and agree with them on $A$.

\emph{A touching ball removes the frontier.}\quad
We can now find a smooth ball touching the unknown part, without
assuming any regularity of the frontier of $A$. Choose $c\in A$
sufficiently close to the center that
\[
 r=\dist(c,V\setminus A)>0,\qquad
 \overline{B(c,r)}\Subset V.
\]
For example, if the center has distance $R$ from $\partial V$,
take $c$ within $R/3$ of it. Then $r<R/3$, and a minimizing
sequence stays in a compact subset of $V$. Thus the distance is
attained at some
\[
 z_*\in\partial B(c,r)\cap(V\setminus A).
\]
On $B(c,r)$ the metrics, the finite harmonic pairs, and the Green
kernels in both variables agree. Smoothness gives equality of
the metric and harmonic-function jets at $z_*$, and the Hessian spanning
condition holds there by our choice of $V$. All the inputs of
Theorem~\ref{local:extension} are therefore satisfied at $z_*$.
It gives equality of the coordinate metrics on a neighborhood of
$z_*$. We may take a small connected ball there meeting $B(c,r)$.
Each paired source difference is now harmonic for the common
metric and vanishes on that intersection. Unique continuation
makes every pair agree on the same ball. The resulting local
isometry thus matches all evaluations.

This isometry agrees with the existing one on every overlap, by
separation of evaluations. Its union with $W$ is an admissible
match whose domain contains $z_*$, contradicting maximality.
Consequently $W=(M_2^e)^\circ$.

\emph{Surjectivity.}\quad
The image is open because $\Phi$ is a local diffeomorphism. It is
also closed in the first interior. Indeed, if $\Phi(p_k)$ tends
to an interior point $q$, a subsequence of $p_k$ converges in
the compact second manifold to $p$. Evaluation continuity again
gives $I_2(p)=I_1(q)$, so $p$ cannot lie on the boundary. Then
$p$ is in the domain and $\Phi(p)=q$. Connectedness of the first
interior makes the image all of it. An injective surjective local
diffeomorphism has a smooth inverse, proving the proposition.
\end{proof}

\subsection{Returning to the original manifold}

The interior isometry still has two duties: it must extend smoothly
through the original boundary, and it must fix every point of the
measured patch, including components where no cap was attached.

\begin{proposition}[Smooth completion and the accessible gauge]
\label{global:completion}
The isometry of Proposition~\ref{global:frontier} extends to a smooth
diffeomorphism $M_2^e\to M_1^e$. It preserves the original copies
of $M$, and its restriction satisfies
\[
 g_2=\Phi^*g_1,\qquad \Phi|_\Gamma=\Id.
\]
\end{proposition}

\begin{proof}
An interior isometry preserves lengths of curves and hence the
intrinsic distances. The metric completion of the interior of a
compact smooth Riemannian manifold with boundary is the manifold
itself with its length metric. To see the only boundary issue,
use a collar to push a boundary-touching curve an arbitrarily
small distance into the interior; its length changes arbitrarily
little. Compactness and local metric comparison then identify
the completion with the original compact manifold.
Thus $\Phi$ and its inverse extend to mutually inverse
homeomorphisms of the two completions. Boundary is mapped to
boundary, since the interiors already correspond bijectively.

These extensions are smooth. In a sufficiently small uniform
boundary collar, the distance $s$ to the boundary is a smooth
function, and its unit gradient is the inward normal flow.
The isometry preserves $s$ and, in the interior, intertwines its
gradients. Fix a small positive level $s=s_0$. The map is already
smooth on that interior hypersurface. Expressing points with
$0\le s\le s_0$ by normal flow from this level expresses $\Phi$
as the composition of this smooth level map with smooth normal
flows on both sides. This extends smoothly to $s=0$. The same
argument applies to the inverse.

Since $\Phi$ is the identity on $E$, continuity makes it the
identity on $\overline E$. The original interiors satisfy
\[
 M_j^\circ=(M_j^e)^\circ\setminus\overline E.
\]
The extended map preserves these complements, and hence their
closures, the original manifolds. Its restriction is therefore a
smooth diffeomorphism $\Phi:M\to M$ with $g_2=\Phi^*g_1$.
It fixes the attached original patch $P_0=\{b>0\}$ pointwise.

To prove the full boundary assertion, take any
$f\in C_c^\infty(\Gamma)$ and let $v_f^j$ now denote the
harmonic functions on the \emph{original} manifolds with boundary
value $f$. On the second manifold both $v_f^2$ and
$v_f^1\circ\Phi$ are $g_2$-harmonic. Their difference has zero
boundary value on $P_0$, where $\Phi$ is the identity.
It also has zero normal derivative there. In fact the equal
energy maps give
\[
 (\partial_{\nu_1}v_f^1)\,dS_{g_1}
 =(\partial_{\nu_2}v_f^2)\,dS_{g_2}
                    \quad\text{on }P_0.
\]
The boundary densities are the same by Lemma~\ref{data:jets},
and the smooth isometry transports the outward unit normal.
Boundary unique continuation, followed by interior continuation,
therefore gives
\[
 v_f^2=v_f^1\circ\Phi\quad\text{throughout }M.
\]
Taking boundary traces yields
\[
 f(q)=f(\Phi(q))\qquad
       (q\in\partial M,\ f\in C_c^\infty(\Gamma)).
\]
For $q\in\Gamma$, these test functions distinguish $q$ from
every other boundary point: choose a function equal to one at
$q$ with support in a small neighborhood avoiding the other
point. Thus $\Phi(q)=q$ for all $q\in\Gamma$.
\end{proof}

\begin{proof}[Proof of Theorem~\ref{main:patch}]
Lemma~\ref{data:jets}, Proposition~\ref{data:cap}, and
Lemma~\ref{data:jet-family} provide the common cap, Green kernels,
and finite harmonic tests. Theorem~\ref{local:extension} applies at
every interior contact
in Proposition~\ref{global:frontier}. Proposition~\ref{global:completion}
gives the required diffeomorphism. All integrations use densities;
the argument required neither an orientation nor connectedness of
the boundary or of $\Gamma$.
\end{proof}
